由引理 \ref{L:faithful realization} 可知，这样定义确实得到一个范畴。自然函子
$$\Sat_G^0\to \on{Mot}_k^*: (X,f)\mapsto \oplus h(X_i)(m_i)$$
是忠实的（但非全的）加法函子。
以下，给定维数为 $\frac{1}{2}(d_{\mmu}+d_{\nu_\bullet})$ 的不可约分支 $Z\subset \Gr_{\mmu}\times_\Gr\Gr_{\nu_\bullet}$，我们用 $[Z]$ 记它在 $\Sat^0_G$ 中诱导的态射。注意，为使用引理 \ref{L:faithful realization}，我们确实需要把 $(\Gr_\mmu,m)$ 的各种「Tate 扭转」也纳入 $\Sat^0_G$ 的对象。

接下来，我们给 $\Sat_G^{0}$ 赋予幺半结构（它与 $\on{Mot}_k^*$ 中通常的张量积结构先验地不同）。
注意对角态射 $\Delta: \Gr_{\mmu}\to \Gr_{\mmu}\times_{\Gr}\Gr_{\mmu}$ 给出 $(\Gr_\mmu,m)$ 在 $\Sat_G^{0}$ 中的恒等态射 $1_{(\Gr_\mmu,m)}$。
我们在连通对象上定义张量积为
\[(\Gr_{\mmu},m)\star (\Gr_{\nu_\bullet},n):=(\Gr_{\mmu,\nu_\bullet},m+n),\]
再按线性性自然延拓到所有对象。
为在态射上定义张量积，只需定义 $1_{(\Gr_{\mmu},m)}\star [Z]$ 与 $[Z']\star 1_{(\Gr_{\nu_\bullet},n)}$，其中 $[Z]\in \Hom_{\Sat_G^{0}}((\Gr_{\nu_\bullet},n'), (\Gr_{\nu'_\bullet},n'))$（相应地 $[Z'] \in \Hom_{\Sat_G^{0}}((\Gr_\mmu,m), (\Gr_{\mu'_\bullet},m'))$）由不可约分支代表。我们定义
$$1_{(\Gr_{\mmu},m)}\star [Z]:=[\Delta\tilde\times Z],\quad [Z']\star 1_{(\Gr_{\nu_\bullet},n)}:=[Z'\tilde\times\Delta],$$
此处用到 $Z$ 是 $K$-不变这一事实，以构造由 \eqref{E-conv Sch} 定义的扭积。

\begin{lem}
\begin{enumerate}
\item 自然（非全）嵌入 $\Sat_G^0\subset\on{Mot}_k^*$ 带有典范的幺半结构，即存在典范同构
\[h(\Gr_{\mmu,\nu_\bullet})(m+n)\cong h(\Gr_{\mmu})(n)\otimes h(\Gr_{\nu_\bullet})(n),\]
它位于 $\on{Mot}^*_k$ 中，并与两个范畴各自的幺半结构相容。
\item $\Sat_G^{0}$ 中存在唯一的交换性约束，使上述幺半函子成为对称幺半函子。
\item 函子 $\Sat_G^0\to\on{Mot}^*_k$ 延拓为对称幺半函子
$$\Sat_G^m\to \on{Mot}_k^*.$$
\end{enumerate}
\end{lem}

由于 $\on{Mot}^*_k$ 幂等完备，而按定义 $\Sat_G^m$ 是 $\Sat_G^0$ 的幂等完备化，最后一条由前两条推出。
把 \cite[\S 2.3, \S 2.4]{Zh1} 中的（等变）上同调换成（等变）Chow 环（参看 \cite{EG}），可以构造映射 $A^*(\Gr_{\mmu})\otimes_{\bQ}A^*(\Gr_{\nu_\bullet})\to A^*(\Gr_{\mmu,\nu_\bullet})$ 与 $A^*(\Gr_{\mmu,\nu_\bullet})\to A^*(\Gr_{\nu_\bullet,\mmu})$，其中 $A^*(-)$ 表示有理 Chow 群。由于所有这些簇都容许胞腔分解，这些映射给出 (1) 与 (2) 中所求结构的候选。
过渡到上同调并利用 \cite[\S 2.3, \S 2.4]{Zh1}，即可看出它们确实是所需的结构。
我们不回顾全部细节，只举一个例子。

\begin{ex}
\label{R: comm via cycle}
设 $G=\GL_2$。
考虑 $\Gr_{\omega_1,\omega_1}$，它分类满足「$\La_i/\La_{i+1}$ 为一维」的格链 $\{\La_2\subset\La_1\subset\La_0=\mO^2\}$。如例 \ref{quasiminuSch} 所述，它同构于 $\bP(\mO_{\bP^1}(-1)\oplus\mO_{\bP^1}(1))$ 的完全化（详细讨论见 \cite[\S B.3]{Zh1}）。回忆其中有一条由条件 $\La_2=\varpi\La_0$ 定义的 $(-2)$-曲线，记作 $\Gr_{\omega_1,\omega_1}^{\omega_2}\subset \Gr_{\omega_1,\omega_1}$。（更确切地说，两种情形下我们都考虑由 $\bP(\mO_{\bP^1}(-1)\oplus\mO_{\bP^1}(1))$ 给出的去完全化；于是对应的约化闭子概形是 $\bP(\mO_{\bP^1}(-1)\oplus\mO_{\bP^1}(1))$ 中的一条 $(-2)$-曲线。）

同构 $h(\Gr_{\omega_1,\omega_1})\cong h(\Gr_{\omega_1})\otimes h(\Gr_{\omega_1})$ 可由轮
$$Z_1+Z_2\subset \Gr_{\omega_1,\omega_1}\times \Gr_{\omega_1}\times \Gr_{\omega_1}$$
实现，其中
$$Z_1=\{(z,\pr_1(z),\pr_1(z))\mid z\in \Gr_{\omega_1,\omega_1}\}, \quad Z_2=\{ (z, \pr_1(z), y)\mid z\in \Gr_{\omega_1,\omega_1}^{\omega_2}, y\in \Gr_{\omega_1}\}.$$
由交换性约束 $c$ 诱导的 $\Gr_{\omega_1,\omega_1}$ 在 $\Sat_G^0$ 中的自同构
\[c':\Gr_{\omega_1,\omega_1}= \Gr_{\omega_1}\star \Gr_{\omega_1}\stackrel{c}{\cong} \Gr_{\omega_1}\star \Gr_{\omega_1}=\Gr_{\omega_1,\omega_1}\]
由
\begin{equation}
\label{E:comm via cycle}
c'=[\Delta]+[\Gr_{\omega_1,\omega_1}^{\omega_2}\times \Gr_{\omega_1,\omega_1}^{\omega_2}]
\end{equation}
给出。这里对角轮 $\Delta$ 与 $\Gr_{\omega_1,\omega_1}^{\omega_2}\times \Gr_{\omega_1,\omega_1}^{\omega_2}$ 是 $\Gr_{\omega_1,\omega_1}\times_{\Gr}\Gr_{\omega_1,\omega_1}$ 的不可约分支，构成
$\Hom_{\Sat_G^{0}}(\Gr_{\omega_1,\omega_1},\Gr_{\omega_1,\omega_1})$ 的一组基。
\end{ex}

下面说明 $\Sat_G^m$ 与通常 Satake 范畴的关系。
设 $\Sat^T_G\subset \on{P}_K(\Gr)$ 为由 $\IC_\mu(i)$（$\mu\in\xcoch(T)/W$，$i\in\bZ$）张成的全子范畴，其中 $\IC_\mu$ 是 $\Gr_\mu$ 上的相交上同调层（它在 Schubert 胞腔上的限制为 $\Ql[d_\mu]$）。由 \cite[引理 5.5.14]{Zh2}（其论证在混合特征下同样成立，无须改动），可以对几何 Satake 作自举，得到如下交换图
\[
\begin{CD}
\Sat^T_G@>\cong>>\on{Rep}_{\Ql}(\hat G^T)\\
@VVV@VVV\\
\Sat_{G\otimes\mO_L}@>\cong>>\on{Rep}_{\Ql}(\hat G).
\end{CD}
\]
其中
\begin{itemize}
\item 群 $\hat{G}^T$ 是 $\hat{G}$ 与 $\bG_m$ 的半直积，$\bG_m$ 对 $\hat G$ 的作用定义如下（细节见 \cite[(5.5.10)]{Zh2}）：记 $\hat G_\ad$ 为 $\hat G$ 的伴随商，它通过共轭作用于 $\hat G$。则复合映射 $\bG_m\to \hat G\to\hat G_\ad$（第一个映射是注 \ref{R: birth of the dual group} (2) 中的余特征标）有平方根 $\rho_\ad:\bG_m\to \hat G_\ad$，后者诱导 $\bG_m$ 在 $\hat G$ 上的一个作用。注意 $\hat G^T$ 带有 Borel 子群 $\hat B^T:=\hat B\rtimes\bG_m$ 与极大环面 $\hat T^T=\hat T\rtimes \bG_m=\hat T\times \bG_m$。
\item 底部的等价即定理 \ref{intro:geomSat} 中的等价。
\item 左侧竖直函子分别是沿 $\Gr\otimes\bar k \to \Gr$ 的拉回函子与沿 $\hat G\to \hat G^T$ 的限制函子。
\end{itemize}

由卷积映射 \eqref{conv m2} 的半小性与分解定理，存在全忠实幺半函子
$$\Sat_G^0\otimes\Ql\to \Sat_G^T,\quad (\Gr_{\mmu},m)\mapsto m_*\Ql_{\Gr_{\mmu}}[d_{\mmu}](m)\in \on{P}_K(\Gr),$$
由 \cite[引理 2.16]{Zh1}，它延拓为等价
$$\Sat^m_G\otimes\Ql\cong \Sat^T_G.$$
因此，按 \cite[(2.4.4)]{Zh1} 对交换性约束作 Koszul 符号改变之后，得到对称幺半范畴的等价
\begin{equation*}
\label{E:geom sat Tate}
\Sat^m_G\otimes\Ql\cong \on{Rep}(\hat{G}^T),
\end{equation*}
我们把它看作动机版的几何 Satake。

\begin{rmk}
\label{R:Tate twist}

如上所述，为了把 $\Sat_G^0$ 定义为 $\on{Mot}^*_k$ 的子范畴，我们确实需要引入 Tate 扭转。因此 $\Sat_G^m\otimes\Ql$ 的 Tannakian 群不是原来的 Langlands 对偶群 $\hat{G}$，而是修正后的 $\hat G^T$。近年来的研究逐渐认识到，在 Langlands 对应的表述中，$\hat G^T$ 或许是更正确的对象（例如见 \cite{BG}）。在某些情形下，可以选取同构 $\hat{G}^T\simeq \hat{G}\times \bG_m$。例如当 $G=\GL_n$ 时，这样的同构可由
\begin{equation}
\label{E: dual group GLn}
\hat{G}^T:=\GL_n\rtimes\bG_m\to \GL_n\times \bG_m,\ (A,t)\mapsto (A\on{diag}\{t^{n-1},t^{n-2},\ldots,1\}, t)
\end{equation}
给出。关于 Langlands 对偶群不同版本的更详细讨论，参看 \cite[\S 5.5]{Zh2}。

该等价诱导 $\bZ[v,v^{-1}]$-代数的同构
$$K(\Sat^m_G)\cong K(\on{Rep}(\hat{G}^T)),$$
其中左端的 $\bZ[v,v^{-1}]$ 由全子范畴 $\{(\Gr_0,m), m\in\bZ\}$ 到 $\Sat_G^m$ 的嵌入诱导，右端的 $\bZ[v,v^{-1}]$ 由全子范畴 $\on{Rep}(\bG_m)\subset \on{Rep}(\hat G^T)$ 的嵌入诱导。于是取 $v=q$ 便得到\emph{典范的} Satake 同构；与 \eqref{Sat isom} 不同，它不依赖于任何选取。当 $\GL_2$ 时，该同构连同 \eqref{E: dual group GLn} 诱导出注 \ref{R: norm Sat} 中的归一化 Satake 同构 $\Sat'^{cl}$。
\end{rmk}

\begin{ex}\label{intersection by geom Sat}
下面是后文应用几何 Satake 的玩具模型。

仍取 $G=\GL_2$。设 $\Gr_{\omega_1,\omega_1^*}$（相应地 $\Gr_{\omega_1^*,\omega_1}$）是分类格链
\begin{equation}\label{Chain 2}
\{\La_2\supset\La_1\subset \La_0=\mO^2\}, \quad \quad \mbox{相应地 } \quad \{\La'_2\subset\La'_1\supset \La_0=\mO_2\}
\end{equation}
的模空间，其中 $\La_0/\La_{1}$ 与 $\La_2/\La_1$（相应地 $\La'_1/\La_0$ 与 $\La'_1/\La'_2$）都是 $k$ 上的一维空间。\zy{[译者按（上下文推断，无外部依据）：此处原文第二条链写作 $\La_0=\mO_2$，据上下文应为 $\La_0=\mO^2$。]}
注意有典范同构
\begin{equation}
\label{E: isomisom1}
\Gr_{\omega_1,\omega_1^*}\cong\Gr_{\omega_1,\omega_1}, \quad \La_2\supset\La_1\subset \La_0\mapsto \varpi\La_2\subset\La_1\subset\La_0,
\end{equation}
与
\begin{equation}
\label{E: isomisom2}
\Gr_{\omega_1^*,\omega_1}\cong\Gr_{\omega_1,\omega_1}, \quad \La'_2\subset\La'_1\supset \La_0\mapsto \varpi\La'_2\subset\varpi\La'_1\subset\La_0,
\end{equation}
故 $\Gr_{\omega_1,\omega_1^*}$ 与 $\Gr_{\omega_1^*,\omega_1}$ 都同构于 $\bP(\mO_{\bP^1}(-1)\oplus\mO_{\bP^1}(1))$ 的完全化。注意
$$\Gr_{\omega_1,\omega_1^*}^0:=\Gr_{\omega_1,\omega_1^*}\times_{\Gr}\Gr_0,\quad\quad \mbox{相应地 }\quad\Gr_{\omega_1^*,\omega_1}^0:=\Gr_{\omega_1^*,\omega_1}\times_{\Gr}\Gr_0$$
是 $\Gr_{\omega_1,\omega_1^*}$（相应地 $\Gr_{\omega_1^*,\omega_1}$）中分类 \eqref{Chain 2} 里满足 $\La_2=\La_0$（相应地 $\La'_2=\La_0$）的格链的闭子概形，因而同构于 $\bP^1$ 的完全化；如例 \ref{quasiminuSch} 所述，它是该曲面上的 $(-2)$-曲线。\zy{[译者按（上下文推断，无外部依据）：此处原文两个下标记作 $\omega_1^*$ 与 $\omega_2$，据上下文应分别为 $\omega_1^*$ 与 $\omega_1$，译文按后者书写。]}

我们考虑 $\Sat_G^{0}$ 中由以下映射复合而成的态射
\[\Gr_0\xrightarrow{\Gr_{\omega_1,\omega_1^*}^0} \Gr_{\omega_1,\omega_1^*}=\Gr_{\omega_1}\star\Gr_{\omega_1^*}\stackrel{c}{\cong}\Gr_{\omega_1^*}\star\Gr_{\omega_1}=\Gr_{\omega_1^*,\omega_1}\xrightarrow{\Gr_{\omega_1^*,\omega_1}^0}\Gr_0,\]
这里把 $\Gr_{\omega_1,\omega_1^*}^0$ 看作 $\Hom(\Gr_0, \Gr_{\omega_1,\omega_1^*})$ 中的元素，把 $\Gr_{\omega_1^*,\omega_1}^0$ 看作 $\Hom(\Gr_{\omega_1^*,\omega_1},\Gr_0)$ 中的元素，$c$ 是交换性约束。结合同构 \eqref{E: isomisom1} 与 \eqref{E: isomisom2}、交换性约束 \eqref{E:comm via cycle}，以及 $\Gr_{\omega_1,\omega_1}^{\omega_2}$ 是 $\Gr_{\omega_1,\omega_1}$ 中 $(-2)$-曲线这一事实，可以看出该复合等于乘 $2$。

另一方面，按几何 Satake，该映射可由 $\GL_2\times \bG_m$ 的有限维表示范畴中的映射
\[\mathbf{1}\to \on{Std}\otimes \on{Std}^*\cong \on{Std}^*\otimes \on{Std} \to \mathbf{1}\]
计算，而后者同样等于乘 $\dim \on{Std}=2$。总之，（卷积）仿射 Grassmann 簇中某些轮之间的相交数，可以借助 Langlands 对偶群的表示论算出。
\end{ex}

最后，我们简要说明如何把上述构造从 $\GL_n$ 推广到 $\mO$ 上的一般约化群 $G$。对一般的约化群，极小余权不够多。例如 $G=E_8$ 时没有非零的极小余权。
故一般情形下，我们把 $\Sat^0_G$ 定义为这样的范畴：对象是 $(\Gr_{\mmu},m)$ 的不交并，其中 $\mmu=(\mu_1,\ldots,\mu_r)$，每个 $\mu_i$ 极小或拟极小（拟极小余权的概念见例如 \cite[\S 2.2.2]{Zh1}）；$(\Gr_{\mmu},m)$ 与 $(\Gr_{\nnu},n)$ 之间的态射由与前面相同的公式给出。

我们需要说明这样确实得到一个范畴（即态射可以复合），并且它有到 $\on{Mot}_k^*$ 的忠实嵌入。
若 $\mu$ 拟极小，$\Gr_\mu$ 不是完全光滑的，而是某个部分旗簇 $G/P$ 上射影锥的完全化；在锥的奇点处爆破即给出奇点的一个显式消解（例如见 \cite[引理 2.12]{Zh1}）。记 $\pi:\wGr_\mu\to \Gr_\mu$ 为这个「消解」，由构造它是 $K$-等变的。不难写出幂等元 $p$，它可以由 $\wGr_\mu\times\wGr_\mu$ 中一个 $K$-不变的轮代表，使得 $(\wGr,p,0)$ 的 étale 实现给出 $\Gr_\mu$ 的相交上同调（由分解定理，它是 $\on{H}^*(\wGr_\mu)$ 的直和项）。
于是有与引理 \ref{L:faithful realization} 类似的引理，只需作显然的修改。重复前面的构造，即得动机 Satake 范畴 $\Sat^m_G$。
